\section{BaseFold From First Principles}

BaseFold combines two ingredients.  We first separate them so the role of each
one is clear.  The first is a concrete IOPP: a protocol realization of the
language-proximity functionality from Figure~\ref{fig:func-prox-language}, with
the language specialized to codewords.  The second is the sumcheck protocol in
Figure~\ref{fig:protocol-sumcheck-randomizes-pcs}: it reduces $f(z)=v$ to a
random evaluation claim.  A polynomial commitment scheme appears only when
these two protocols are run together with the same challenges.

This section follows the BaseFold construction and terminology~[BaseFold23].
When it discusses IOPPs and proximity, it relies on the proof-model vocabulary
introduced above from RVW13 and BCS16~[RVW13,BCS16].

The central design question is simple.  Sumcheck wants to fold the message
representation: split the current message as $m=(m_L\mid m_R)$ and replace it
by $m_L+\alpha m_R$ for a verifier challenge $\alpha$.  But the verifier does
not have the raw message.  It has a commitment to an encoded oracle
$\pi=\Enc(m)$.  BaseFold is the answer to the question:
\[
  \text{can we fold the committed encoding as if we had folded the message?}
\]
The code is built so that the answer is yes.

\subsection{BaseFold As A Proximity Instantiation}

The proximity functionality was defined generically in
Figure~\ref{fig:func-prox-language}.  BaseFold instantiates that functionality
with a code language.  If \(C\subseteq\F^N\) is the public code at one layer,
then the language is simply
\[
  L_C=C=\{\Enc(m):m\in\F^K\}\subseteq\F^N.
\]
Thus the ideal target is
\[
  \mathcal F_{\mathrm{Prox}}^{L_C,\delta}.
\]
A concrete BaseFold IOPP realizes this code-language proximity target, up to
its soundness error, by testing that the committed word is close to
\emph{some} valid codeword.  This statement is only about proximity to
\(L_C\).  It does not say which polynomial was encoded, and it does not by
itself prove \(f(z)=v\).  The PCS will use this proximity statement as one
component.

\subsection{Concrete Foldable Codes}

Now specialize to the concrete foldable-code structure used by the BaseFold
paper.  For $i=0,\ldots,d$, let $C_i$ be a linear code with message tensor
domain $\B^{k_i}$ and codeword tensor domain $\B^{n_i}$.  The actual message
length is $K_i=2^{k_i}$, and the actual codeword length is $N_i=2^{n_i}$.
At each fold level one Boolean coordinate is removed, so
\[
  k_{i+1}=k_i+1,
  \qquad
  n_{i+1}=n_i+1.
\]
Let
\[
  \Enc_i:\F^{\B^{k_i}}\to\F^{\B^{n_i}}
\]
be its encoder.  Equivalently, with row-vector convention,
\[
  \Enc_i(m)=mG_i,
\]
where $G_i$ is the public generator matrix whose columns are indexed by
$\B^{n_i}$.

The code family is foldable because the next generator matrix is built from the
previous one using a public random diagonal matrix $T_i$ whose rows and columns
are indexed by $\B^{n_i}$.  Diagonal means that
\[
  T_i[a,a']=0
  \qquad\text{whenever } a\ne a'.
\]
In the concrete random foldable code, the diagonal entries
\[
  \{T_i[a,a]\}_{a\in\B^{n_i}}
\]
are sampled independently and uniformly from $\F$, then fixed as part of the
public code description:
\[
  G_{i+1}
  =
  \begin{bmatrix}
    G_i & G_i \\
    G_iT_i & G_i(T_i+I_i)
  \end{bmatrix}.
\]
Here $I_i$ is the identity matrix on $\F^{\B^{n_i}}$.
The two block columns are indexed by $(a,0)$ and $(a,1)$ for
$a\in\B^{n_i}$.  The diagonal entries of $T_i$ are part of the code
description.  Write
\[
  t_{i,a}=T_i[a,a]\in\F.
\]
Then the two interpolation points for coordinate $a$ are $t_{i,a}$ and
$t_{i,a}+1$.  These points are not chosen by the prover and are not verifier
challenges; they are fixed by the concrete random foldable code.

Split a message $m\in\F^{\B^{k_{i+1}}}$ as
\[
  m=(m_L\mid m_R),
  \qquad m_L,m_R\in\F^{\B^{k_i}}.
\]
If $\pi_{i+1}=\Enc_{i+1}(m)$, then for every $a\in\B^{n_i}$,
\[
  \pi_{i+1}[a,0]
  =
  \Enc_i(m_L)[a]+t_{i,a}\Enc_i(m_R)[a],
\]
\[
  \pi_{i+1}[a,1]
  =
  \Enc_i(m_L)[a]+(t_{i,a}+1)\Enc_i(m_R)[a].
\]
Thus these two symbols are evaluations of the same line
\[
  L_a(X)=\Enc_i(m_L)[a]+X\Enc_i(m_R)[a]
\]
at the two public code points $t_{i,a}$ and $t_{i,a}+1$.  If the verifier
later samples a challenge $\alpha_i\in\F$, evaluating the same line at
$\alpha_i$ gives
\[
  L_a(\alpha_i)
  =
  \Enc_i(m_L+\alpha_i m_R)[a].
\]
This is the concrete BaseFold fold.  It is a local interpolation rule on
codeword symbols, and its correctness comes from the displayed generator-matrix
decomposition.

Now we can name the two sides of the same operation.  The message fold is
\[
  \MsgFold_{\alpha_i}(m)=m_L+\alpha_i m_R\in\F^{\B^{k_i}}.
\]
The codeword fold is the local interpolation rule
\[
  \CodeFold_{\alpha_i}(\pi_{i+1})[a]=L_a(\alpha_i),
  \qquad a\in\B^{n_i},
\]
where $L_a$ is reconstructed from the two codeword symbols at the
code-specified points $t_{i,a}$ and $t_{i,a}+1$.  Foldability is the
commuting statement
\[
  \CodeFold_{\alpha_i}(\Enc_{i+1}(m))
  =
  \Enc_i(\MsgFold_{\alpha_i}(m)).
\]
This equation is the intuition that the protocol relies on: sumcheck asks for
the left-to-right message fold, while the verifier can only test the
top-to-bottom codeword fold through committed oracle openings.

\subsection{Minimum Distance And Johnson Radius}

The proof needs a way to talk about how close words can be to codewords, and
how close different codewords can be to each other.  The first tool is minimum
distance.  Half the minimum distance is the uniqueness radius: inside that
radius, a received word can be close to at most one codeword.  That is useful,
but it is not the whole story.  Just outside the uniqueness radius, there may be
more than one nearby codeword.  The proof also needs to know that there are not
too many of them.  Johnson's bound is the standard tool for controlling that
nearby list.

Let $C\subseteq\F^{\B^n}$ be a code.  For words $x,c\in\F^{\B^n}$, write
\[
  \Delta(x,c)
  =
  \frac{|\{a\in\B^n:x[a]\ne c[a]\}|}{2^n}
\]
for ordinary relative Hamming distance.  The relative minimum distance of $C$
is
\[
  \Delta_C
  =
  \min_{\substack{c,c'\in C\\ c\ne c'}}\Delta(c,c').
\]
Thus two distinct codewords disagree on at least a $\Delta_C$ fraction of
coordinates.

\paragraph{Random foldable codes have good minimum distance.}
For BaseFold's random foldable code (RFC), the minimum distance is not assumed by
magic; it is proved from the random diagonal matrices.  Let
$N_i=2^{n_i}$ be the actual codeword length at layer $i$.  The BaseFold
minimum-distance theorem says, in this paper's notation, that the random choice
of the foldable code gives an explicit lower bound
\[
  \Delta_{C_d}\ge \rho_{\mathrm{RFC}}
\]
except with small probability over the sampled diagonal matrices.  The quantity
$\rho_{\mathrm{RFC}}\in[0,1]$ is computed from the code parameters
$(k_0,d,c,|\F|)$; Appendix C of BaseFold gives the recurrence, and the theorem's
failure probability is, for example, on the order of $d\cdot 2^{-\lambda}$ for
security parameter $\lambda$.

The proof style is important.  Because the code is linear, the minimum distance
is the minimum Hamming weight of a nonzero codeword.  So the proof studies zeros:
how many coordinates of $\Enc_i(m)$ can vanish for a nonzero message $m$?
Inductively assume that level $i$ has few zeros for every nonzero message.  At
the next level, split $m=(m_L\mid m_R)$ and write the two previous-layer
encodings as
\[
  A=\Enc_i(m_L),
  \qquad
  B=\Enc_i(m_R).
\]
For each previous-layer coordinate $a\in\B^{n_i}$, the next codeword contains
two line evaluations whose slope is controlled by the random diagonal entry
$t_{i,a}$.  If $A[a]$ and $B[a]$ are not both zero, then the event that one of
those new symbols vanishes is essentially one linear equation in the random
field element $t_{i,a}$, hence has probability about $1/|\F|$.

A naive proof would now union bound over all nonzero messages, but that loses too
much: there are too many messages.  BaseFold's proof refines the union bound by
partitioning messages according to the common zero set
\[
  S=\{a:A[a]=0\text{ and }B[a]=0\}.
\]
Large sets $S$ are dangerous because those coordinates are automatically zero in
both halves, but rank-nullity says there are few messages that can have a large
common zero set.  Small sets $S$ allow many messages, but then many independent
random diagonal entries would need to hit unlucky roots.  The proof balances
these two effects: count messages by the size of $S$, apply a binomial-style
tail bound for the fresh random diagonal entries, and union bound over all
possible $S$.  That is how the random foldable code gets a concrete lower bound
on $\Delta_{C_d}$.

\paragraph{Minimum distance gives unique decoding below half distance.}
If a word $v\in\F^{\B^n}$ satisfies
\[
  \Delta(v,c)<\Delta_C/2,
  \qquad
  \Delta(v,c')<\Delta_C/2
\]
for two codewords $c,c'\in C$, then the triangle inequality gives
\[
  \Delta(c,c')
  \le
  \Delta(c,v)+\Delta(v,c')
  <
  \Delta_C.
\]
This contradicts the definition of $\Delta_C$ unless $c=c'$.  So inside radius
$\Delta_C/2$, there is at most one nearby codeword.  This is the basic
minimum-distance argument.

\paragraph{Johnson controls lists beyond half distance.}
Unique decoding says there is at most one nearby codeword below
$\Delta_C/2$.  List decoding asks a softer question beyond that radius: even if
there are several nearby codewords, can we bound how many there are?  Johnson's
bound gives such a radius using only the minimum distance of the code.

For a slack parameter $\gamma\in(0,1]$, define the Johnson-radius function
\[
  J_\gamma(\lambda)
  =
  1-\sqrt{1-\lambda(1-\gamma)}
  \qquad\text{for }\lambda\in[0,1].
\]
This is a slightly shrunken version of the usual Johnson radius
$1-\sqrt{1-\lambda}$.  The shrinkage gives the proof room for strict
inequalities and union bounds.

The proof intuition is a counting tension.  If many codewords are all within
relative distance $\rho$ of the same received word $v$, then each of them agrees
with $v$ on about a $1-\rho$ fraction of coordinates.  Averaging over pairs of
nearby codewords, many of those agreements must occur on the same coordinates,
so the nearby codewords agree with each other too often.  But distinct codewords
are required to disagree on at least a $\Delta_C$ fraction of coordinates.
Solving this tension gives a controlled list size as long as
$\rho$ is below the Johnson radius.

BaseFold uses this through the following one-step no-free-repair lemma.  Assume
the code minimum distances are monotone in the BaseFold direction, so the
top-layer minimum distance $\Delta_{C_d}$ is a valid lower bound for the layers
below.
Set
\[
  R_{\mathrm{BF}}
  =
  J_\gamma(J_\gamma(\Delta_{C_d})).
\]

\begin{lemma}[One-step no-free-repair lemma]
\label{lem:one-step-no-free-repair}
For any layer $i=0,\ldots,d-1$, any word
$v\in\F^{\B^{n_{i+1}}}$, and any $\eta\le R_{\mathrm{BF}}$, if
\[
  \Delta^*(v,C_{i+1})>\eta,
\]
then
\[
  \Pr_{\alpha_i\leftarrow\F}
  \left[
    \Delta(\CodeFold_{\alpha_i}(v),C_i)
    \le
    \eta-\gamma
  \right]
  \le
  \frac{2}{\gamma^3|\F|}.
\]
Here $\Delta^*$ is the paired/coset distance used by BaseFold: a parent pair is
counted as bad if either of its two entries disagrees with the candidate
codeword.
\end{lemma}

The proof idea matches the foldable-code structure above.  Any parent word
$v$ can be uniquely written as two point evaluations of a coordinatewise line,
so the possible folds $\CodeFold_{\alpha_i}(v)$ trace the affine family
$u+\alpha_i u'$.  If too many challenges $\alpha_i$ made this fold close to
some child codeword of $C_i$, Johnson/list decoding would force those nearby
child codewords to share enough structure to identify two codewords
$w,w'\in C_i$ agreeing with $u,u'$ on many coordinates.  Foldability then
combines $w$ and $w'$ into a parent codeword in $C_{i+1}$ that is close to
$v$ in paired/coset distance, contradicting
$\Delta^*(v,C_{i+1})>\eta$.  This is the formal reason adaptive free repairs
are rare.

\subsection{BaseFold IOPP}

Recall the split from Figure~\ref{fig:func-prox-language}: the semantic target
is the ideal functionality \(\mathcal F_{\mathrm{Prox}}^{L,\delta}\), while an
IOPP is the protocol that realizes that target with sublinear openings.  In
this subsection the language is the top code language
\[
  L_{C_d}=C_d=\{\Enc_d(m):m\in\F^{\B^{k_d}}\}.
\]
The BaseFold IOPP does not prove an evaluation claim such as \(f(z)=v\).  It is
the concrete protocol that realizes proximity to \(L_{C_d}\): an accepting
transcript should mean that the committed top oracle is close to some codeword
of \(C_d\), except with the protocol's soundness error.  BaseFold does this by
asking the prover to commit to folded oracle layers and then auditing a small
number of local fold relations between adjacent layers.

First read the protocol in the honest direction.  Fix a number of fold layers
$d\in\mathbb{N}$.  The public objects are the field $\F$, the Boolean coordinate
alphabet $\B$, the log message and codeword dimensions
$k_i,n_i\in\mathbb{N}$ for $i=0,\ldots,d$, with
$k_{i+1}=k_i+1$ and $n_{i+1}=n_i+1$, the encoders
\[
  \Enc_i:\F^{\B^{k_i}}\to\F^{\B^{n_i}},
  \qquad
  C_i=\Enc_i(\F^{\B^{k_i}}),
\]
the code-specified interpolation points $t_{i,a}\in\F$ for
$i=0,\ldots,d-1$ and $a\in\B^{n_i}$, and a query count
$Q_{\mathrm{BF}}\in\mathbb{N}$.  The IOPP input oracle is a committed tensor
$\pi_d\in\F^{\B^{n_d}}$.  In an honest execution, the prover's private input is
a message tensor $m_d\in\F^{\B^{k_d}}$, and the committed top oracle is
\[
  \pi_d=\Enc_d(m_d)\in\F^{\B^{n_d}}.
\]
For each layer $i=d-1,\ldots,0$, the verifier samples a fresh challenge
$\alpha_i\leftarrow\F$, sends it to the prover, and then the prover commits to
the next child oracle.  In
the concrete coefficient-style BaseFold convention of the previous subsection,
the honest prover splits $m_{i+1}=(m_L\mid m_R)$ with
$m_L,m_R\in\F^{\B^{k_i}}$ and sets
\[
  m_i=\MsgFold_{\alpha_i}(m_{i+1})=m_L+\alpha_i m_R.
\]
The honest child oracle is then
\[
  \pi_i=\Enc_i(m_i).
\]
Equivalently, and this is the form the verifier can test, for every
$a\in\B^{n_i}$ the value $\pi_i[a]$ is obtained by interpolating the line through
the two parent symbols
\[
  (t_{i,a},\pi_{i+1}[a,0])
  \qquad\text{and}\qquad
  (t_{i,a}+1,\pi_{i+1}[a,1])
\]
and evaluating that line at $\alpha_i$.  Thus the honest invariant is
\[
  \pi_i=\Enc_i(m_i),
\]
where $m_i$ is the folded message determined by the challenges already sampled.
The verifier cannot check this invariant everywhere, so
Figure~\ref{fig:protocol-basefold-iopp} authenticates the committed oracles and
then checks a few complete parent-child paths.  Notice the protocol box itself
does not force the prover to compute the honest child oracle before committing
to it; it lets the prover commit a claimed child oracle.  The query phase is
where the verifier tests whether those claimed children are consistent with the
folds of their parents.

\begin{protocolbox}{Protocol: BaseFold IOPP}{fig:protocol-basefold-iopp}
\label{fig:protocol-basefold-iopp-query}
\noindent Public parameters: \(d,Q_{\mathrm{BF}}\in\mathbb{N}\), codes
\(\mathcal{C}_i\subseteq\F^{\B^{n_i}}\), interpolation points
\(t_{i,a}\in\F\) for \(i=0,\ldots,d-1\) and \(a\in\B^{n_i}\).

\smallskip
\begin{enumerate}[leftmargin=0.24in,label=\textbf{P\arabic*.},itemsep=0.08em,topsep=0.1em,parsep=0pt]
  \item \(\prover\): \(R_d\gets\mathcal F_{\mathrm{MsgOracle}}.\mathsf{Commit}(\pi_d)\).
  \item \(\prover\to\verifier\): \(R_d\).
  \item For \(i=d-1,d-2,\ldots,0\):
        \begin{enumerate}[leftmargin=0.28in,label=\textbf{P\arabic*.\alph*.},itemsep=0.05em,topsep=0.05em,parsep=0pt]
          \item \(\verifier\): \(\alpha_i\leftarrow\F\).
          \item \(\verifier\to\prover\): \(\alpha_i\).
          \item \(\prover\): choose a claimed child oracle
                \(\pi_i\in\F^{\B^{n_i}}\).
          \item \(\prover\): \(R_i\gets\mathcal F_{\mathrm{MsgOracle}}.\mathsf{Commit}(\pi_i)\).
          \item \(\prover\to\verifier\): \(R_i\).
        \end{enumerate}
  \item For \(q=0,1,\ldots,Q_{\mathrm{BF}}-1\):
        \begin{enumerate}[leftmargin=0.28in,label=\textbf{P\arabic*.\alph*.},itemsep=0.05em,topsep=0.05em,parsep=0pt]
          \item \(\verifier\): \(a_{q,d-1}\leftarrow\B^{n_{d-1}}\).
          \item \(\verifier\to\prover\): \(a_{q,d-1}\).
          \item \(\verifier\): for every \(i=0,\ldots,d-1\), set
                \(a_{q,i}\in\B^{n_i}\) to the length-\(n_i\) prefix of
                \(a_{q,d-1}\).
          \item For \(i=d-1,d-2,\ldots,0\):
                \begin{enumerate}[leftmargin=0.32in,label=\textbf{P\arabic*.\alph*.\roman*.},itemsep=0.04em,topsep=0.04em,parsep=0pt]
                  \item \(\prover,\verifier\):
                        \(x_0\gets\mathcal F_{\mathrm{MsgOracle}}.\mathsf{Open}(R_{i+1},(a_{q,i},0))\).
                  \item \(\prover,\verifier\):
                        \(x_1\gets\mathcal F_{\mathrm{MsgOracle}}.\mathsf{Open}(R_{i+1},(a_{q,i},1))\).
                  \item \(\prover,\verifier\):
                        \(x_\star\gets\mathcal F_{\mathrm{MsgOracle}}.\mathsf{Open}(R_i,a_{q,i})\).
                  \item \(\verifier\): if \(x_0=\bot\) or \(x_1=\bot\) or
                        \(x_\star=\bot\), return reject.
                  \item \(\verifier\): let \(L(X)\in\F[X]\) be the line with
                        \(L(t_{i,a_{q,i}})=x_0\) and
                        \(L(t_{i,a_{q,i}}+1)=x_1\).
                  \item \(\verifier\): if \(L(\alpha_i)\ne x_\star\), return reject.
                \end{enumerate}
        \end{enumerate}
  \item \(\verifier\): for every \(a\in\B^{n_0}\), set
        \(x_a\gets\mathcal F_{\mathrm{MsgOracle}}.\mathsf{Open}(R_0,a)\).
  \item \(\verifier\): if \(x_a=\bot\) for any \(a\in\B^{n_0}\), return reject.
  \item \(\verifier\): set \(\widehat\pi_0\gets(x_a)_{a\in\B^{n_0}}\).
  \item \(\verifier\): if \(\widehat\pi_0\notin\mathcal{C}_0\), return reject.
  \item \(\verifier\): return accept.
\end{enumerate}
\end{protocolbox}

\paragraph{Soundness view.}
Now read Figure~\ref{fig:protocol-basefold-iopp} as a test against a cheating
prover.  In the ideal message-oracle model of Figure~\ref{fig:func-msg-oracle},
the verifier sees handles and opened coordinates, while an extractor can read
the full tensors that the prover committed to the functionality.  In
particular, after an accepting execution the extractor can inspect the committed
top oracle
$\pi_d\in\F^{\B^{n_d}}$.

We want to rule out the case where $\pi_d$ is far from the top code $C_d$.
For a non-bottom layer $i$ and a codeword $c\in C_i$, let
$\Delta^*(v,c)$ denote the BaseFold paired/coset relative distance between
$v\in\F^{\B^{n_i}}$ and $c$: group coordinates as $(a,0),(a,1)$ for
$a\in\B^{n_i-1}$ and count the fraction of pairs on which $v$ and $c$ disagree
in at least one entry.  Then define the distance from a word to the code by
\[
  \Delta^*(\pi_d,C_d)
  =
  \min_{c\in C_d}\Delta^*(\pi_d,c).
\]
Choose a nearest codeword
\[
  c_d^\star
  \in
  \operatorname*{arg\,min}_{c\in C_d}\Delta^*(\pi_d,c),
\]
and write
\[
  c_d^\star=\Enc_d(m_d^\star),
  \qquad
  m_d^\star\in\F^{\B^{k_d}},
  \qquad
  e_d=\pi_d-c_d^\star\in\F^{\B^{n_d}}.
\]
We assume the encoders are injective, so a codeword identifies its message.

There are two useful ways to measure the residual $e_d$.  The familiar one is
ordinary relative Hamming weight:
\[
  |e_d|_{\mathrm{Ham}}
  =
  \frac{
    \bigl|\{x\in\B^{n_d}:e_d[x]\ne0\}\bigr|
  }{
    2^{n_d}
  }.
\]
BaseFold's proof uses the coarser paired support, because one local fold check
reads a pair of parent coordinates.  Group the top-layer coordinates as
$(a,0),(a,1)$ for $a\in\B^{n_d-1}$ and define
\[
  |e_d|_{\mathrm{pair}}
  =
  \frac{
    \bigl|\{a\in\B^{n_d-1}:e_d[a,0]\ne0\text{ or }e_d[a,1]\ne0\}\bigr|
  }{
    2^{n_d-1}
  }.
\]
With $c_d^\star$ chosen by the same paired/coset distance, the mental picture is
\[
  \Delta^*(\pi_d,C_d)
  =
  \Delta^*(\pi_d,c_d^\star)
  =
  |e_d|_{\mathrm{pair}},
\]
with the normalization above.  Let
$\delta\in[0,1]$ denote the effective top-layer proximity distance in the
BaseFold theorem: this is the above paired/coset distance capped by the
Johnson/list-decoding radius $R_{\mathrm{BF}}$ defined above.
Thus $\delta=0$ means the committed top oracle is already a codeword.  Larger
$\delta$ means that every codeword is farther away, so more parent-pair positions
would have to be repaired before the oracle could look valid.  The proof does
not need this number to keep growing all the way to $1$: beyond the
code-dependent radius, the theorem treats the word as already far enough.
For scale, if the top layer has $1024$ parent pairs and the nearest codeword
differs from $\pi_d$ in at least one entry of $300$ of those pairs, then
$\Delta^*(\pi_d,C_d)=300/1024\approx0.29$ under this normalization.  That is a
large distance: the word is not a codeword with a few typos; almost a third of
the pair positions need to be explained away.

\paragraph{Now switch to soundness.}
Presume the bad case: $\delta$ is large, but the prover still wants the verifier
to accept.  The bottom layer is checked directly in
Figure~\ref{fig:protocol-basefold-iopp}, so an accepting prover needs the
committed bottom oracle $\pi_0$ to be a valid word in the public base code
$C_0$.  Thus a far top oracle must somehow become a valid bottom oracle along
the fold chain.

The concrete game at one layer is simple.  The parent oracle $\pi_{i+1}$ is
already committed.  Then the verifier samples $\alpha_i$.  At that moment the
parent commitment determines a unique truthful folded word
\[
  w_i=\CodeFold_{\alpha_i}(\pi_{i+1})\in\F^{\B^{n_i}}.
\]
The prover now commits to the child oracle $\pi_i$.  At each child coordinate
$a\in\B^{n_i}$, it can either set $\pi_i[a]=w_i[a]$, which is locally truthful,
or set $\pi_i[a]\ne w_i[a]$, which creates a local inconsistency that a query
path can catch.  With this notation, there are only two ways the distance to the
code can drop as the protocol moves from layer $i+1$ to layer $i$.

\paragraph{Free repairs.}
It helps to start with a weaker, non-adaptive game.  Imagine that, before seeing
$\alpha_i$, the prover has already picked a target child codeword
\[
  \widehat c_i\in C_i.
\]
The prover is still allowed to lie in the child oracle.  The simplification is
only that the desired target and the intended lie positions are fixed before the
challenge is known.

For a child coordinate $a\in\B^{n_i}$, the committed parent pair
\[
  \pi_{i+1}[a,0],
  \qquad
  \pi_{i+1}[a,1]
\]
determine a line $L_a(X)$ through the two public code points
$t_{i,a}$ and $t_{i,a}+1$.  The honest folded child value at this coordinate is
\[
  L_a(\alpha_i).
\]
In the fixed-target game, a free repair at $a$ means that the truthful value
$L_a(\alpha_i)$ happens to hit the preselected target value $\widehat c_i[a]$.
That is a single linear condition in the random challenge $\alpha_i$.  This is
the simple local picture: if the target is fixed before the challenge,
accidental repairs are easy to believe rare.

If $L_a(\alpha_i)\ne\widehat c_i[a]$ and the prover still sets
$\pi_i[a]=\widehat c_i[a]$, then the repair is not free.  It is a local lie
against the forced fold value $w_i[a]$, and the query phase may catch it.

The real prover is stronger.  It does not have to name $\widehat c_i$ before
seeing $\alpha_i$.  The adaptive free-repair event is that $\pi_{i+1}$ is far
from $C_{i+1}$, but the forced truthful fold $w_i$ is close to some convenient
codeword of $C_i$.  That convenient codeword may depend on $\alpha_i$.  So the
bad event is not
\[
  w_i[a]=\widehat c_i[a]
\]
for a fixed $\widehat c_i$.  It is the word-level statement that $w_i$ has
become close to the code $C_i$ at all.

Lemma~\ref{lem:one-step-no-free-repair} is the replacement for the fixed-target
local argument.  The bad event it excludes is that, for some layer $i$,
\[
  \Delta(w_i,C_i)
  \le
  \min\{\Delta^*(\pi_{i+1},C_{i+1}),R_{\mathrm{BF}}\}-\gamma,
\]
where $\Delta(x,C_i)=\min_{c\in C_i}\Delta(x,c)$ and $\Delta(x,c)$ is ordinary
relative Hamming distance on the child layer.  The parameter $\gamma>0$ is the
slack used by the proximity analysis.  The lemma says this bad event has small
probability over the verifier's random fold challenges.
Condition on the complementary good-challenge event.  From now on, random folds
do not give the prover a noticeable global repair for free.

\paragraph{Paid repairs.}
Under the good-challenge event, the truthful folded word $w_i$ cannot be much
closer to $C_i$ than the parent distance allows.  If the prover still wants a
cleaner child layer, it must commit to a child oracle $\pi_i$ that differs from
$w_i$.  These differences are the prover's lying mass at layer $i$: they may
move $\pi_i$ closer to the code, but they create local parent-child fold
violations.  Define
\[
  \beta_i
  =
  \Delta(\pi_i,w_i),
\]
where $\Delta$ is ordinary relative Hamming distance on the child coordinate
domain $\B^{n_i}$.  Thus $\beta_i$ is the fraction of child coordinates where the
committed child disagrees with the actual fold of the committed parent.  In
plain language, $\beta_i$ is how much the prover lies at layer $i$.

Equivalently, call a node $(i,a)$ bad when the committed child value
$\pi_i[a]$ is inconsistent with the two parent values
$\pi_{i+1}[a,0]$ and $\pi_{i+1}[a,1]$ under the fold challenge $\alpha_i$.
Then $\beta_i$ is the fraction of bad nodes in layer $i$.

\paragraph{Telescoping.}
Now keep a distance ledger across the layers.  The top oracle starts far from
the top code: its effective distance is $\delta$.  The bottom oracle must end as
an actual codeword of $C_0$, because Figure~\ref{fig:protocol-basefold-iopp}
checks $C_0$ membership directly.  So an accepting prover has to make the
distance-to-code drop from $\delta$ at the top to $0$ at the bottom.

The good-challenge event says that honest folding cannot explain much of this
drop for free.  Therefore, if the distance drops from layer $i+1$ to layer $i$,
the drop must mostly be paid for by local disagreements between the committed
child $\pi_i$ and the forced folded word $w_i$.  The telescoping argument proves
that the total lying mass $\sum_i\beta_i$ has to be large.  The rest of this
paragraph is just notation for that ledger.

The cap in the ledger is
\[
  R_{\mathrm{BF}}
  =
  J_\gamma(J_\gamma(\Delta_{C_d})).
\]
This is the Johnson/list-decoding radius from
Lemma~\ref{lem:one-step-no-free-repair}; beyond this radius the proof simply
treats a word as far enough.  For each layer, define the capped distance
\[
  \delta^{(i)}
  =
  \min\{R_{\mathrm{BF}},\Delta^*(\pi_i,C_i)\}
\]
for non-bottom layers.  For notational convenience at the base layer, read
$\Delta^*(\pi_0,C_0)$ as ordinary relative distance $\Delta(\pi_0,C_0)$, and set
$\delta^{(0)}=0$ when the prover passes the exhaustive check $\pi_0\in C_0$.
Otherwise the verifier rejects immediately.  With this notation,
$\delta^{(d)}=\delta$ and $\delta^{(0)}=0$.

The good-challenge event gives
\[
  \Delta(w_i,C_i)>\delta^{(i+1)}-\gamma.
\]
In words: the truthful fold of the parent, $w_i$, is still almost as far from
$C_i$ as the parent layer was from $C_{i+1}$, up to the slack $\gamma$.
On the other hand, the triangle inequality gives
\[
  \Delta(w_i,C_i)
  \le
  \Delta(w_i,\pi_i)+\Delta^*(\pi_i,C_i)
  =
  \beta_i+\Delta^*(\pi_i,C_i).
\]
In words: if $w_i$ is far from $C_i$ but the committed child $\pi_i$ is closer,
then $w_i$ and $\pi_i$ must differ in many positions.  Those differences are
exactly the layer-$i$ lying mass measured by $\beta_i$.
If $\delta^{(i)}\ge \delta^{(i+1)}-\gamma$, then the desired lower bound on
$\beta_i$ is already trivial.  Otherwise the cap is not active at layer $i$, so
$\delta^{(i)}=\Delta^*(\pi_i,C_i)$, and the two displayed inequalities imply
\[
  \beta_i \ge \delta^{(i+1)}-\delta^{(i)}-\gamma
\]
up to the harmless strict/weak inequality convention in the paper.  Summing this
per-layer lying lower bound over the layers gives the core bound
\[
  \sum_{i=0}^{d-1}\beta_i \ge \delta-\gamma d,
\]
because the distance starts at $\delta^{(d)}=\delta$ and ends at
$\delta^{(0)}=0$.

The verifier's paths are attempts to catch this lying mass.  For query repetition
$q$, Figure~\ref{fig:protocol-basefold-iopp} samples
$a_{q,d-1}\in\B^{n_{d-1}}$ and follows its prefixes down the chain.  At level
$i$, the verifier checks the child value $\pi_i[a_{q,i}]$ against the fold of
the two parent values $\pi_{i+1}[a_{q,i},0]$ and
$\pi_{i+1}[a_{q,i},1]$.  Thus an accepting cheating prover needs every sampled
path to avoid the parent-child positions counted by the $\beta_i$ terms.

The path matters because the lies may be spread across layers.  If the
verifier first chose a random layer $i$ and then checked a random coordinate in
that layer, the catch probability would be an average,
\[
  \frac{1}{d}\sum_{i=0}^{d-1}\beta_i.
\]
That would lose a factor of $d$.  A path check is different: one sampled top
coordinate induces one child coordinate at every lower layer, so the same query
gets a chance to hit lies at all layers.  Intuitively, the path asks
whether this whole ancestor chain is consistent, not whether one isolated layer
looks clean.

There is one more technical step in the paper.  After fixing good fold
challenges, BaseFold gives the prover a cleanup for free.  Along a path, once
the prover has lied at some node, there is no reason to create extra downstream
lies merely to propagate that same changed value.  Those extra lies only make the
path easier to catch.  The proof therefore modifies the internal oracle layers
$\pi_{d-1},\ldots,\pi_1$ without changing the endpoints $\pi_d$ or $\pi_0$ and
without increasing the rejection probability.  This cleanup keeps only the
necessary first lies on paths and makes all avoidable downstream values
consistent with them.

After this normalization, one queried path rejects exactly when it hits at least
one remaining bad node.  In that normalized view, the one-path rejection
probability is
\[
  \sum_{i=0}^{d-1}\beta_i,
\]
and therefore at least $\delta-\gamma d$.  This is where the correlated
top-to-bottom path check enters: after cleanup, $\sum_i\beta_i$ is the mass of
unavoidable first lies that paths must cross, not a double-count of messy
downstream aftershocks.

With $Q_{\mathrm{BF}}$ independently sampled paths, the probability that a word
at effective relative distance $\delta$ survives is therefore bounded by
\[
  (1-\delta+\gamma d)^{Q_{\mathrm{BF}}}
  \le
  \exp\!\left(-Q_{\mathrm{BF}}(\delta-\gamma d)\right).
\]

This gives an information-theoretic extraction story for the concrete
proximity realization.  The extractor reads $\pi_d$ from
$\mathcal F_{\mathrm{MsgOracle}}$ and, conceptually, finds a nearest codeword such
as $c_d^\star$.  Below the unique-decoding radius, the corresponding message
$m_d^\star$ is unique.  The brute-force search may be inefficient; an efficient
extractor would need an efficient decoder for this code family.

That decoder caveat matters.  Codes with strong efficient decoders are not the
main object in the random foldable-code regime.  If extraction were forced to
rely on a decoder with guaranteed relative radius
$\rho_{\mathrm{dec}}\in[0,1]$ smaller than the proximity threshold
\(\delta\), then the effective threshold for an efficiently extractable theorem would be
$\rho_{\mathrm{dec}}$.  To recover a target soundness level $\lambda>0$ at that
smaller threshold, the verifier would need more query paths, roughly scaling
like $Q_{\mathrm{BF}}\approx\lambda/\rho_{\mathrm{dec}}$.  That restores
soundness, but increases proof size, commitment openings, and verifier work.

For this subsection, the important endpoint is only proximity:
\[
  \pi_d \approx \Enc_d(m_d^\star).
\]
It is not yet the PCS statement that the polynomial represented by
$m_d^\star$ evaluates to $v$ at $z$.  The next subsection adds sumcheck back in
to bind an evaluation claim to the same encoded message.

\subsection{BaseFold PCS: Sumcheck Superimposed With IOPP}

We can now build the PCS evaluation protocol.  The public claim is
\[
  f(z)=v,
  \qquad z\in\F^d,\quad v\in\F,
\]
where the prover has committed to an encoding of the polynomial data for
$f$.  Sumcheck alone can randomize this claim, but it does not make the committed
object smaller.  The BaseFold proximity protocol alone can fold a committed
codeword down to a tiny codeword, but its semantic target is only
\(\mathcal F_{\mathrm{Prox}}^{L_C,\delta}\); it does not say which evaluation
claim should hold.  The PCS is obtained by doing both operations with the same
verifier challenges.

This is the point of the superposition.  Sumcheck changes the large claim
$f(z)=v$ into a terminal random claim $f(r)=u$.  At the same time, the
BaseFold proximity protocol changes the committed object itself.  It folds the
top encoding of $f$ into a bottom encoding of a smaller polynomial, then a
smaller one again, until the last encoded object is just one scalar.  Thus the
final randomized sumcheck claim is not checked against the large top-level
object.  It is checked against the tiny bottom encoding, which the verifier can
open and inspect directly.

For this concrete presentation, use the simplifying BaseFold setting $k_0=0$,
so $k_i=i$, $K_0=2^{k_0}=1$, and the bottom code $C_0$ encodes one scalar.  To
match the BaseFold paper, let $y_d\in\F^{\B^d}$ be the coefficient tensor of the
multilinear polynomial $f_d=f$.  That is,
\[
  f_d(X_0,\ldots,X_{d-1})
  =
  \sum_{b\in\B^d} y_d[b]\prod_{i=0}^{d-1}X_i^{b_i}.
\]
The top oracle is supposed to be
\[
  \pi_d=\Enc_d(y_d).
\]
Here $y_d$ is being serialized as the message of the code.  It is not a Boolean
evaluation tensor in this paragraph, so we do not write $f_d=\MLE(y_d)$.

There are now two ledgers to track.  The first ledger is the sumcheck ledger.
With $\chi_z$ as defined earlier, set
\[
  H(Y)=f_d(Y)\chi_z(Y).
\]
For Boolean $b\in\B^d$, $\chi_z(b)=\chi_b(z)$, so
\[
  f_d(z)
  =
  \sum_{b\in\B^d}f_d(b)\chi_z(b)
  =
  \sum_{b\in\B^d}H(b).
\]
The starting claim $f_d(z)=v$ is therefore the Boolean-sum claim
$\sum_{b\in\B^d}H(b)=v$.

After the verifier has sampled the later challenges
$r_{i+1},\ldots,r_{d-1}$, define
\[
  H_{i+1}(Y_0,\ldots,Y_i)
  =
  H(Y_0,\ldots,Y_i,r_{i+1},\ldots,r_{d-1}).
\]
At the start of round $i$, sumcheck maintains the claim
\[
  s_{i+1}
  =
  \sum_{b\in\B^{i+1}}H_{i+1}(b).
\]
The prover sends a univariate polynomial claiming to be
\[
  g_i(X)=\sum_{b\in\B^i}H_{i+1}(b,X).
\]
The check $g_i(0)+g_i(1)=s_{i+1}$ verifies the boundary between the old sum and
the new one.  Then the verifier samples $r_i$ and the maintained sumcheck claim
becomes
\[
  s_i=g_i(r_i)
  =
  \sum_{b\in\B^i}H_i(b),
  \qquad
  H_i(Y_0,\ldots,Y_{i-1})
  =
  H_{i+1}(Y_0,\ldots,Y_{i-1},r_i).
\]

The second ledger is the committed-representation ledger.  It starts from a
tensor $y_d\in\F^{\B^d}$ and its codeword $\pi_d=\Enc_d(y_d)$, but the tensor
does not have a meaning until we choose a representation.  The tag
$\tau\in\{\mathsf{coeff},\mathsf{eval}\}$ says how to read it:
\[
  f_d(X)=
  \begin{cases}
    \displaystyle\sum_{b\in\B^d}y_d[b]\prod_{t=0}^{d-1}X_t^{b_t},
      & \tau=\mathsf{coeff},\\[0.9em]
    \MLE(y_d)(X),
      & \tau=\mathsf{eval}.
  \end{cases}
\]
Suppose layer $i+1$ is honest and encodes
$y_{i+1}\in\F^{\B^{i+1}}$, which represents the current polynomial
$f_{i+1}:\F^{i+1}\to\F$ according to the same tag $\tau$.  Split
\[
  y_{i+1}=(y_L\mid y_R),
  \qquad y_L,y_R\in\F^{\B^i}.
\]
The same challenge $r_i$ used by sumcheck defines the folded tensor
\[
  y_i[b]=
  \begin{cases}
    y_L[b]+r_i y_R[b],
      & \tau=\mathsf{coeff},\\
    (1-r_i)y_L[b]+r_i y_R[b],
      & \tau=\mathsf{eval},
  \end{cases}
  \qquad b\in\B^i,
\]
and in either representation this tensor represents
\[
  f_i(X_0,\ldots,X_{i-1})
  =
  f_{i+1}(X_0,\ldots,X_{i-1},r_i).
\]
The coefficient formula is specialization of the last variable.  The
evaluation formula is the ordinary MLE fold on the last Boolean coordinate.
These are different arrays, but they represent the same restricted polynomial
when they are read in their own bases.

The honest child oracle is
\[
  \pi_i=\Enc_i(y_i).
\]
The encoders and public interpolation points are part of the same typed
instantiation: for the chosen tag $\tau$, the local code fold must produce an
encoding of the folded tensor above.  The verifier cannot read the whole oracle
chain, so the BaseFold protocol checks sampled parent-child interpolation
relations.  Those local checks are exactly the BaseFold mechanism for realizing
proximity to a valid folded-oracle chain, namely the sequence
\[
  \Enc_d(y_d),\Enc_{d-1}(y_{d-1}),\ldots,\Enc_0(y_0).
\]

Now the two ledgers meet.  Sumcheck folds $H=f_d\chi_z$ down to the scalar
\[
  s_0=H(r_0,\ldots,r_{d-1})
      =f_d(r_0,\ldots,r_{d-1})\chi_z(r).
\]
The committed-polynomial ledger folds $f_d$ down to
\[
  f_0=f_d(r_0,\ldots,r_{d-1}),
\]
and, in the $k_0=0$ setting, $y_0=f_0$ is one field element.  If
$\chi_z(r)\ne0$, the verifier computes
\[
  u=s_0/\chi_z(r)
\]
and checks that the opened bottom oracle is $\Enc_0(u)$.  This is the small
clear check that replaces the original large committed evaluation claim.

\begin{protocolbox}{Protocol: BaseFold PCS Evaluation}{fig:protocol-basefold-pcs-eval}
\scriptsize
\begin{description}[leftmargin=0.2in,style=nextline,itemsep=0.16em,topsep=0.2em,parsep=0pt]
  \item[Public parameters.]
        $\tau\in\{\mathsf{coeff},\mathsf{eval}\}$;
        $\Enc_i:\F^{\B^{k_i}}\to\F^{\B^{n_i}}$ for $i=0,\ldots,d$;
        $k_{i+1}=k_i+1$ and $n_{i+1}=n_i+1$;
        $t_{i,a}\in\F$ for $i=0,\ldots,d-1$ and $a\in\B^{n_i}$;
        $Q_{\mathrm{BF}}\in\mathbb{N}$.
  \item[Prover input.]
        $(y_d,z,v)$, where $y_d\in\F^{\B^d}$, $z\in\F^d$, and $v\in\F$.
  \item[Verifier input.]
        $(z,v)$.
  \item[Interpret $y_d$.]
        \(\prover\):
        \[
          f_d(X)\gets
          \begin{cases}
            \displaystyle\sum_{b\in\B^d}y_d[b]\prod_{t=0}^{d-1}X_t^{b_t},
              & \tau=\mathsf{coeff},\\[0.9em]
            \MLE(y_d)(X),
              & \tau=\mathsf{eval}.
          \end{cases}
        \]
  \item[Commit phase.]
        \begin{enumerate}[leftmargin=0.28in,label=\textbf{C\arabic*.},itemsep=0.08em,topsep=0.1em,parsep=0pt]
          \item \(\prover\): \(\pi_d\gets\Enc_d(y_d)\in\F^{\B^{n_d}}\).
          \item \(\prover\): \(C_d\gets\mathcal F_{\mathrm{MsgOracle}}.\mathsf{Commit}(\pi_d)\).
          \item \(\prover\to\verifier\): \(C_d\).
        \end{enumerate}
  \item[Open phase.]
        \begin{enumerate}[leftmargin=0.28in,label=\textbf{O\arabic*.},itemsep=0.08em,topsep=0.1em,parsep=0pt]
          \item \(\prover\): \(H(Y)\gets f_d(Y)\chi_z(Y)\).
          \item \(\verifier\): \(s_d\gets v\).
          \item For $i=d-1,d-2,\ldots,0$:
            \begin{enumerate}[leftmargin=0.24in,label=\textbf{\alph*.},itemsep=0.06em,topsep=0.06em,parsep=0pt]
              \item \(\prover\):
                \[
                  g_i(X)\gets
                  \sum_{b\in\B^i}H(b,X,r_{i+1},\ldots,r_{d-1}).
                \]
                \(\prover\to\verifier\): \(g_i\).
              \item \(\verifier\): if \(\deg(g_i)>2\) or
                \(g_i(0)+g_i(1)\ne s_{i+1}\), return reject.
              \item \(\verifier\): \(r_i\leftarrow\F\).\\
                \(\verifier\to\prover\): \(r_i\).\\
                \(\verifier\): \(s_i\gets g_i(r_i)\).
              \item \(\prover\): split \(y_{i+1}=(y_L\mid y_R)\) with
                \(y_L,y_R\in\F^{\B^i}\), and set
                \[
                  y_i[b]\gets
                  \begin{cases}
                    y_L[b]+r_i y_R[b],
                      & \tau=\mathsf{coeff},\\
                    (1-r_i)y_L[b]+r_i y_R[b],
                      & \tau=\mathsf{eval},
                  \end{cases}
                  \qquad b\in\B^i.
                \]
              \item \(\prover\): for every \(a\in\B^{n_i}\), compute
                \[
                  L_{i,a}(X)\gets
                  \operatorname{line}\bigl((t_{i,a},\pi_{i+1}[a,0]),
                  (t_{i,a}+1,\pi_{i+1}[a,1])\bigr),
                  \qquad
                  \pi_i[a]\gets L_{i,a}(r_i).
                \]
              \item \(\prover\): \(C_i\gets\mathcal F_{\mathrm{MsgOracle}}.\mathsf{Commit}(\pi_i)\).\\
                \(\prover\to\verifier\): \(C_i\).
            \end{enumerate}
          \item \(\verifier\): \(r\gets(r_0,\ldots,r_{d-1})\).
          \item \(\verifier\): if \(\chi_z(r)=0\), return reject.
          \item \(\verifier\): \(u\gets s_0/\chi_z(r)\).
          \item \(\prover,\verifier\): for every \(a\in\B^{n_0}\), set
            \(o_a\gets\mathcal F_{\mathrm{MsgOracle}}.\mathsf{Open}(C_0,a)\).
          \item \(\verifier\): if \(o_a=\bot\) for any \(a\in\B^{n_0}\), return reject.
          \item \(\verifier\): set \(\widehat\pi_0\gets(o_a)_{a\in\B^{n_0}}\).
          \item \(\verifier\): if \(\widehat\pi_0\ne\Enc_0(u)\), return reject.
          \item For $q=0,\ldots,Q_{\mathrm{BF}}-1$:
            \begin{enumerate}[leftmargin=0.24in,label=\textbf{\alph*.},itemsep=0.06em,topsep=0.06em,parsep=0pt]
              \item \(\verifier\): \(a_{q,d-1}\leftarrow\B^{n_{d-1}}\).\\
                \(\verifier\to\prover\): \(a_{q,d-1}\).
              \item \(\verifier\): for every \(i=0,\ldots,d-1\), set
                \(a_{q,i}\in\B^{n_i}\) to the length-\(n_i\) prefix of
                \(a_{q,d-1}\).
              \item For $i=d-1,d-2,\ldots,0$:
                \begin{enumerate}[leftmargin=0.24in,label=\textbf{\roman*.},itemsep=0.06em,topsep=0.06em,parsep=0pt]
                  \item \(\prover,\verifier\): set
                    \(x_0\gets\mathcal F_{\mathrm{MsgOracle}}.\mathsf{Open}(C_{i+1},(a_{q,i},0))\),
                    \(x_1\gets\mathcal F_{\mathrm{MsgOracle}}.\mathsf{Open}(C_{i+1},(a_{q,i},1))\),
                    and \(x_\star\gets\mathcal F_{\mathrm{MsgOracle}}.\mathsf{Open}(C_i,a_{q,i})\).
                  \item \(\verifier\): if \(x_0=\bot\) or \(x_1=\bot\) or \(x_\star=\bot\), return reject.
                  \item \(\verifier\):
                    \[
                      L(X)\gets
                      \operatorname{line}\bigl((t_{i,a_{q,i}},x_0),
                      (t_{i,a_{q,i}}+1,x_1)\bigr);
                    \]
                  \item \(\verifier\): if \(L(r_i)\ne x_\star\), return reject.
                \end{enumerate}
            \end{enumerate}
          \item \(\verifier\): return accept.
        \end{enumerate}
\end{description}
\end{protocolbox}

Figure~\ref{fig:protocol-basefold-pcs-eval} is the compact protocol form of the
walkthrough above.  The PCS claim itself is representation-free: the verifier is
trying to check $f_d(z)=v$.  The protocol box chooses a representation only for
the committed witness tensor $y_d$ and the folded tensors that follow it.  For
$\tau=\mathsf{coeff}$, $y_i$ stores multilinear coefficients, and folding is
coefficient specialization.  For $\tau=\mathsf{eval}$, $y_i$ stores Boolean
evaluations, and folding is the affine MLE fold on the last Boolean coordinate.
The verifier should not mix these two readings; in particular,
$f_i=\MLE(y_i)$ is true only in the evaluation case.

The sumcheck equations determine the terminal claim $f_d(r)=u$.  The BaseFold
path checks and proximity analysis do not say that every verifier-facing oracle
is error-free.  They say that, except with the backend soundness error, the
committed oracle chain is close to a valid folded-codeword chain for typed
tensors $y_d,y_{d-1},\ldots,y_0$, with bottom value $y_0=u$.  The same field
element $r_i$ is used both as the sumcheck challenge in round $i$ and as the
code-fold challenge from layer $i+1$ to layer $i$.  This is why the two
transcripts meet at the same object $(r,u)$ instead of proving two unrelated
facts.

The soundness picture is just these two invariants stacked on top of each other.
Sumcheck contributes the algebraic invariant: except with the usual
Schwartz--Zippel error from the random challenges, a prover who passes the
round checks has translated the original claim $f_d(z)=v$ into the terminal
claim $f_d(r)=u$.  The BaseFold proximity protocol contributes the commitment
invariant: its realization of
\(\mathcal F_{\mathrm{Prox}}^{L_C,\delta}\), plus the fold checks, says that a
prover who passes has committed to a top word close to an actual codeword, and
that the same challenges are consistent with a nearby folded-codeword chain
ending at the opened scalar $u$.  The code is what makes this compression
robust.  The verifier does not trust a few local openings as a codeword proof by
themselves; it relies on the BaseFold proximity analysis to turn those
correlated local checks into the statement that the committed object is close to
one extracted polynomial, whose folded endpoint is $u$.

\subsection{Knowledge Extraction}

\paragraph{Framing.}
Soundness says that a false evaluation claim is rejected with high probability.
Knowledge soundness asks for more.  If a prover can convince the verifier with
non-negligible probability, then an extractor should recover the nearby
polynomial enforced by the commitment and proximity proof.  The apparent
obstacle is exactly the one we identified for the standalone proximity
realization: from the top oracle $\pi_d$ alone,
recovery would look like a decoding problem for the code $C_d$, and this code
family is not assumed to have an efficient decoder in the relevant parameter
regime.

The PCS extractor avoids decoding.  It uses the prover as an evaluation oracle
for that extracted nearby polynomial.  Each accepting execution with challenge vector
$r\in\F^d$ yields a terminal value $u$, and the soundness argument above says
that, except with negligible probability,
\[
  u=f_d(r),
\]
where $f_d$ is the unique polynomial whose encoding is close to the committed
top oracle.  Enough such evaluations determine the coefficient tensor of $f_d$.

\paragraph{Required invariants.}
The extraction argument relies on three facts already established by the pieces
above.

First, the proximity analysis gives a unique nearby polynomial.  If the
prover passes with non-negligible probability, then the top oracle is close to
some codeword $\Enc_d(y_d)$:
\[
  \Delta^*(\pi_d,\Enc_d(y_d))\le\delta.
\]
For parameters below the uniqueness radius, this $y_d$ is unique.  Conceptually,
the accepted transcript identifies one nearby polynomial, even if the extractor
does not know how to decode $\pi_d$ to find it directly.

Second, the superimposed protocol makes accepted endpoints correct.  Fix an
accepting run with challenge vector $r=(r_0,\ldots,r_{d-1})$.  Sumcheck says that
the original claim has been translated to the terminal claim $f_d(r)=u$.  The
BaseFold proximity protocol says that the committed oracle chain folds, under
the same $r_i$, close to a valid folded-codeword chain ending at $u$.  If the
prover tried to make the accepted endpoint come from a different folded
polynomial, then some layer would have to switch from being close to the
extracted folded codeword to being close to another codeword.  The local fold
consistency, the code distance, and the proximity bound rule this out except
with negligible probability.

Third, random evaluation points span the coefficient space.  For
$\rho=(\rho_0,\ldots,\rho_{d-1})\in\F^d$, define the expansion vector
\[
  \operatorname{exp}_d(\rho)\in\F^{\B^d},
  \qquad
  \operatorname{exp}_d(\rho)[b]
  =
  \prod_{t=0}^{d-1}\rho_t^{b_t}.
\]
If $y_d$ is the coefficient tensor of $f_d$, then
\[
  f_d(\rho)=\sum_{b\in\B^d}y_d[b]\operatorname{exp}_d(\rho)[b].
\]
Thus $2^d$ evaluations at points whose expansion vectors are linearly
independent give a square linear system for the $2^d$ coefficients of $y_d$.
Random points have this independence property with high probability; the proof
uses the Schwartz--Zippel lemma to show that a fresh random point almost always
adds a new independent row until the coefficient space is full.

The proof packages this last sampling step as a predicate-forking lemma.  In
this setting, the predicate is ``the expansion vectors collected so far are
linearly independent.''  If the prover accepts with probability $\varepsilon$ on
a random challenge vector, the lemma lets an expected polynomial-time extractor,
with overhead proportional to $1/\varepsilon$, find enough accepting challenge
vectors that also satisfy this independence predicate.

\begin{protocolbox}{Extractor: BaseFold PCS Knowledge Extraction}{fig:extractor-basefold-pcs}
\footnotesize
\begin{description}[leftmargin=0.22in,style=nextline,itemsep=0.18em,topsep=0.2em,parsep=0pt]
  \item[Input.]
        Black-box access to a prover $P^\star$ that convinces the verifier in
        Figure~\ref{fig:protocol-basefold-pcs-eval} with non-negligible
        probability $\varepsilon$, and access to the oracle strings committed to
        $\mathcal F_{\mathrm{MsgOracle}}$.
  \item[Goal.]
        Recover a coefficient tensor $y_d\in\F^{\B^d}$ such that
        $\Delta^*(\pi_d,\Enc_d(y_d))\le\delta$ and the represented polynomial
        satisfies $f_d(z)=v$.
  \item[Collect successful challenges.]
        Repeatedly run $P^\star$ with fresh challenge vectors
        $\rho\leftarrow\F^d$, using $\rho_i$ as the round-$i$ sumcheck and
        code-fold challenge.  Keep only accepting executions.  A standard
        predicate-forking argument lets the extractor obtain $2^d$ accepting
        challenge vectors
        \[
          \rho^{(1)},\ldots,\rho^{(2^d)}
        \]
        whose expansion vectors
        $\operatorname{exp}_d(\rho^{(j)})$ are linearly independent.
  \item[Record endpoint values.]
        For each accepted run $j$, read the terminal sumcheck value
        $s_0^{(j)}$ and compute
        \[
          u_j=s_0^{(j)}/\chi_z(\rho^{(j)}),
        \]
        in the simplified presentation where $\chi_z(\rho^{(j)})\ne0$.  By the
        soundness invariant, except with negligible probability,
        $u_j=f_d(\rho^{(j)})$.
  \item[Interpolate.]
        Solve the linear system
        \[
          \sum_{b\in\B^d}
            y_d[b]\operatorname{exp}_d(\rho^{(j)})[b]
          =
          u_j,
          \qquad j=1,\ldots,2^d.
        \]
        The rows are linearly independent, so the solution is unique.  Output
        the recovered coefficient tensor $y_d$.
\end{description}
\end{protocolbox}

\paragraph{Why this works.}
The extractor is efficient because it turns knowledge extraction into
interpolation, not decoding.  The code is still essential, but its job is
different: it prevents the prover from answering each challenge with a different
polynomial.  Once the BaseFold proximity and sumcheck soundness arguments ensure
that every accepted endpoint is an evaluation of the same extracted nearby
polynomial $f_d$, the extractor can learn $f_d$ from evaluations.  This is why
the full PCS can have an
efficient knowledge extractor even though the standalone proximity protocol by
itself only gave a conceptual nearest-codeword extractor.

\subsection{Systematic Codes And External Input Commitments}

The concrete BaseFold paper presentation above assumes that BaseFold owns one
commitment to the whole top oracle.  A systematic presentation separates the
logical codeword from the way its coordinates are authenticated.  It also uses a
modified code: the top-level code is not the generic RFC codeword with a wrapper
around some coordinates.  It is a systematic code whose message block is literal
and whose parity block is built with the same foldable recurrence.

Fix a depth $d$.  At level $i$, let the message length be $K_i=2^{k_i}$ and
the parity length be $M_i=2^{n_i}$.  The parity encoder is a fold-compatible RFC
encoder
\[
  E_i:\F^{\B^{k_i}}\to\F^{\B^{n_i}}
\]
defined recursively as follows.  Start with the public base parity encoder
$E_0$.  For each $i=0,\ldots,d-1$, the code description contains a public
diagonal schedule $T_i$ whose rows and columns are indexed by $\B^{n_i}$.
Write
\[
  t_{i,a}=T_i[a,a]\in\F
  \qquad(a\in\B^{n_i}).
\]
Given
\[
  y=(y_L\mid y_R)\in\F^{\B^{k_{i+1}}},
  \qquad
  y_L,y_R\in\F^{\B^{k_i}},
\]
let
\[
  L=E_i(y_L)\in\F^{\B^{n_i}},
  \qquad
  R=E_i(y_R)\in\F^{\B^{n_i}}.
\]
Then $E_{i+1}(y)\in\F^{\B^{n_i}\times\B}$ is defined by
\[
  E_{i+1}(y)[a,0]
  =
  L[a]+t_{i,a}\bigl(R[a]-L[a]\bigr),
\]
\[
  E_{i+1}(y)[a,1]
  =
  L[a]+(t_{i,a}+1)\bigl(R[a]-L[a]\bigr)
\]
for every $a\in\B^{n_i}$.  Thus each parity pair is the RFC line between the
two child parity encodings $L$ and $R$, evaluated at the code-specified points
$t_{i,a}$ and $t_{i,a}+1$.

The systematic codeword length is
\[
  N_i=K_i+M_i.
\]
The message coordinates are indexed by $\B^{k_i}$, and the parity coordinates
are indexed by $\B^{n_i}$.  The actual systematic code is
\[
  C_i^{\mathrm{sys}}(y)
  =
  \bigl(y\mid E_i(y)\bigr)
  \in
  \F^{\B^{k_i}}\times\F^{\B^{n_i}}.
\]
Equivalently, with row-vector convention,
\[
  E_i(y)=yG_i^{\mathrm{par}},
\]
and
\[
  G_i^{\mathrm{sys}}
  =
  \begin{bmatrix}
    I_{K_i} & G_i^{\mathrm{par}}
  \end{bmatrix},
\]
so the first $K_i$ coordinates are the message itself and the last $M_i$
coordinates are parity.  This is not the generic RFC generator matrix with some
columns declared systematic.  It is a new systematic generator matrix obtained
by adjoining an identity block to the RFC parity generator.  The distance proof
for this modified systematic code is omitted here and handled as a separate
code-theoretic result.

The point of using this RFC encoder as the parity block is that adjacent
systematic encoders still fold correctly.  Split
$y\in\F^{\B^{k_{i+1}}}$ as
\[
  y=(y_L\mid y_R),
  \qquad
  y_L,y_R\in\F^{\B^{k_i}}.
\]
For the raw systematic block, the two parent entries
$y_L[j]$ and $y_R[j]$ are values at the finite points $0$ and $1$.  Folding at
the verifier challenge $\theta\in\F$ gives
\[
  y^\star[j]=(1-\theta)y_L[j]+\theta y_R[j].
\]
For the parity block, the two parent symbols above $a\in\B^{n_i}$ are the two
values defined by the recursive encoder above.  They are evaluations of the same
line $L[a]+X(R[a]-L[a])$ at the code-specified points $t_{i,a}$ and
$t_{i,a}+1$.

Use two instances of $\mathcal F_{\mathrm{MsgOracle}}$ from
Figure~\ref{fig:func-msg-oracle} at the top layer $d$.  The first instance is
external and has a public handle $R_{\mathrm{sys}}$.  Its implementation is
opaque to BaseFold: it may be a virtual message-oracle handle whose opening
protocol authenticates some larger object and derives the logical value $y[j]$.  The
second instance is owned by the backend and has handle $R_{\mathrm{par}}$ for
the parity vector $\widetilde y=E_d(y)$.  Thus a top-layer opening is dispatched
by coordinate type:
\[
  c[j]=
  \begin{cases}
    y[j], & 0\le j<K_d,\\
    \widetilde y[j-K_d], & K_d\le j<N_d.
  \end{cases}
\]
If $0\le j<K_d$, the verifier checks the opening against $R_{\mathrm{sys}}$.
If $K_d\le j<N_d$, it checks the opening against $R_{\mathrm{par}}$.
After the top layer, the folded oracles
$\pi_{d-1},\ldots,\pi_0$ are backend-owned objects; they can be committed with
ordinary Merkle implementations of $\mathcal F_{\mathrm{MsgOracle}}$.

The fold equations use the two coordinate blocks differently.  On systematic
coordinates, $\CodeFold$ specializes to $\MsgFold$ for the representation stored
there.

If those entries are Boolean evaluations, the fold with verifier challenge
$\theta\in\F$ is the MLE fold
\[
  y^\star[a]=(1-\theta)y[(a\mid0)]+\theta y[(a\mid1)],
  \qquad
  \MLE(y^\star)(r)=\MLE(y)(r,\theta).
\]
If those entries are coefficients, the same affine fold is coefficient
specialization:
\[
  f^\star(X)=f(X,\theta).
\]
On parity coordinates, the verifier cannot use the raw message fold because
parity symbols are not message entries.  It uses the code-dependent public
points $t_{i,a}$ and $t_{i,a}+1$ supplied by the RFC.  Interpolating
the two opened parity symbols at $\theta$ gives
\[
  L[a]+\theta\bigl(R[a]-L[a]\bigr)
  =
  E_i\bigl((1-\theta)y_L+\theta y_R\bigr)[a],
\]
where the last equality uses linearity of $E_i$.  This is why the local codeword
view generalizes the raw message fold:
\[
  \CodeFold_\theta(C_{i+1}^{\mathrm{sys}}(y))
  =
  C_i^{\mathrm{sys}}(\MsgFold_\theta(y)).
\]
On systematic symbols this equation is the direct message fold.  On parity
symbols it is enforced by the encoder's structured local relation.  The parity
commitment is what gives the externally authenticated message entries enough
redundancy for the proximity argument.
